\section{Selector Dasar: Elemen, Class, dan Id}

Selector adalah mekanisme CSS untuk menentukan elemen HTML mana yang akan menerima gaya tertentu. Pemilihan selector yang tepat adalah keterampilan fundamental karena menentukan jangkauan dan specificity aturan CSS. CSS menyediakan berbagai jenis selector---dari yang paling sederhana (elemen tunggal) hingga yang sangat spesifik (kombinasi bertingkat). Pada section ini, fokus pada tiga selector dasar yang paling sering digunakan: selector elemen, class, dan id \cite{mdnCss}.

\textbf{Selector elemen} (type selector) menargetkan semua elemen HTML dengan nama tag tertentu. Misalnya, \code{p \{ color: navy; \}} akan mengubah warna teks semua elemen \code{<p>} di halaman. Selector ini memiliki specificity rendah dan cocok untuk mengatur gaya dasar (base style) yang berlaku secara global. \textbf{Selector class} menggunakan tanda titik (\code{.}) diikuti nama class, misalnya \code{.card}. Sebuah class dapat diterapkan pada banyak elemen berbeda, menjadikannya selector paling serbaguna dan paling direkomendasikan untuk styling komponen modular \cite{webdevCss}.

\textbf{Selector id} menggunakan tanda pagar (\code{\#}) diikuti nama id, misalnya \code{\#header}. Karena id harus unik per halaman, selector ini hanya menargetkan satu elemen. Specificity id sangat tinggi, yang membuatnya sulit di-override. Oleh karena itu, praktik modern merekomendasikan penggunaan class daripada id untuk styling; id lebih cocok digunakan untuk anchor link dan atribut JavaScript (\code{getElementById}) \cite{cssTricks}.

\begin{table}[htbp]
\centering
\begin{tabular}{|P{2.5cm}|P{2.5cm}|P{3cm}|P{4cm}|}
\hline
\textbf{Jenis Selector} & \textbf{Sintaks} & \textbf{Specificity} & \textbf{Kapan Digunakan} \\
\hline
Elemen & \code{p}, \code{h1} & 0-0-0-1 & Base style untuk seluruh tag \\
\hline
Class & \code{.card} & 0-0-1-0 & Komponen modular; paling direkomendasikan \\
\hline
Id & \code{\#header} & 0-1-0-0 & Elemen unik; hindari untuk styling \\
\hline
Universal & \code{*} & 0-0-0-0 & Reset global \\
\hline
Gabungan & \code{p.highlight} & Jumlah komponen & Target lebih spesifik \\
\hline
Descendant & \code{nav a} & Jumlah komponen & Elemen di dalam konteks \\
\hline
Child & \code{ul > li} & Jumlah komponen & Anak langsung saja \\
\hline
\end{tabular}
\caption{Jenis-jenis selector CSS dan tingkat specificity-nya}
\end{table}

Selain tiga selector dasar, CSS mendukung \textbf{selector gabungan} untuk targeting yang lebih presisi. \textbf{Descendant selector} (\code{nav a}) menargetkan semua \code{<a>} di dalam \code{<nav>}, termasuk yang bersarang dalam. \textbf{Child selector} (\code{ul > li}) hanya menargetkan anak langsung. \textbf{Grouping selector} (\code{h1, h2, h3}) menerapkan gaya yang sama ke beberapa selector sekaligus. \textbf{Adjacent sibling} (\code{h2 + p}) menargetkan elemen yang langsung mengikuti elemen lain. Penggunaan selector gabungan yang tepat mengurangi kebutuhan menambah class berlebihan di HTML \cite{mdnCss}.

\begin{webcode}[language=CSS, caption={Selector elemen, class, id, dan gabungan}]
/* Selector elemen: semua paragraf */
p { line-height: 1.6; }

/* Selector class: komponen card */
.card {
  padding: 16px;
  border: 1px solid #ddd;
  border-radius: 8px;
}

/* Selector id: elemen unik */
#hero-banner { background-color: #2a6fb0; }

/* Gabungan: paragraf dengan class 'highlight' */
p.highlight { background-color: yellow; }

/* Descendant: tautan di dalam nav */
nav a { text-decoration: none; }

/* Child: hanya li langsung di bawah ul */
ul.menu > li { display: inline-block; }

/* Grouping: heading yang sama */
h1, h2, h3 { font-family: 'Segoe UI', sans-serif; }
\end{webcode}

\begin{catatan}
\textbf{Class vs Id untuk Styling.} Meskipun id memiliki specificity lebih tinggi, praktik modern merekomendasikan penggunaan class untuk hampir semua styling. Alasan: (1) class dapat dipakai berulang, id tidak; (2) specificity id terlalu tinggi---sulit di-override tanpa \code{!important}; (3) banyak framework CSS (Bootstrap, Tailwind) dibangun sepenuhnya berbasis class. Gunakan id untuk atribut \code{for} pada label, anchor navigation (\code{href="\#section"}), dan JavaScript \cite{cssTricks}.
\end{catatan}

